\chapter{Divisibility and the Euclidean Algorithm}\label{ch:divisibility}

\begin{center}\itshape
The oldest non-trivial algorithm still in daily use is Euclid's, from around 300 BC.\\
Its worst case is governed by the Fibonacci numbers, as Lam\'e proved in 1844, in an early systematic analysis of an algorithm's complexity.
\end{center}


A calculation with two integers will open the door. Divide $252$ by $198$:
\[
252=198+54,\qquad 198=3\cdot54+36,\qquad
54=36+18,\qquad36=2\cdot18.
\]
The last nonzero remainder is $18$. It divides both starting numbers, and every common divisor of $252$ and $198$ divides it. That is why the greatest common divisor is $18$.

The calculation raises two questions that matter more than the answer. Why does replacing a pair by the smaller number and its remainder keep every common divisor? And why must the remainders eventually stop? The first question is answered by subtraction. The second is answered by the fact that a strictly decreasing list of nonnegative integers cannot continue forever. Neither answer relies on guessing the gcd.

A gcd is also a linear combination. Reading the divisions backwards gives $18=54-36=4\cdot54-198=4\cdot252-5\cdot198$. Those coefficients are not decoration: they will become modular inverses in Chapter~\ref{ch:congruences}. Euclid's algorithm is both a decision procedure and a machine for constructing witnesses.

\begin{think}
In the calculation above, why is the last nonzero remainder a \emph{greatest} common divisor, rather than merely a divisor shared by the two numbers? Answer using divisibility, not numerical size.
\end{think}


\section{Division and ideals of \texorpdfstring{$\Z$}{Z}}

% t1-pass2 motive 7efe951b
What can division do that subtraction cannot? If $a=bq+r$, the remainder $r$ is smaller than $b$ and keeps exactly the information we need about common divisors. We will prove that the remainder exists and is unique, including when $a$ is negative. Then we will collect all integer combinations $ax+by$ into one set. Its least positive element is the gcd, and writing it as a combination gives B\'ezout's identity. The abstract word \emph{ideal} comes later; here the object is something we can calculate with.

\begin{reminder}[What is allowed in a remainder?]
Here $\N=\{0,1,2,\ldots\}$ and $\N^*=\{1,2,\ldots\}$. We write $a\mid b$ if $b=ak$ for some integer $k$; thus every integer divides $0$, but $0$ divides only $0$. A prime is a positive integer greater than $1$ whose positive divisors are $1$ and itself. The number $1$ is a unit, not a prime: allowing it as a prime would give $6=2\cdot3=1\cdot2\cdot3=1^2\cdot2\cdot3$, and uniqueness would never get out of the starting gate.

In Euclidean division the remainder is always non-negative, even for a negative dividend. Thus $-17=5(-4)+3$, rather than $5(-3)-2$. The quotient is $\lfloor a/b\rfloor$, the greatest integer at most $a/b$.
\end{reminder}

\begin{theorem}[Euclidean division]\label{thm:div}
For $a\in\Z$ and $b\in\N^*$ there is a unique pair $(q,r)\in\Z^2$ with $a=bq+r$ and $0\le r<b$.
\end{theorem}
\begin{proof}
\emph{Existence.} Consider all nonnegative integers of the form
$a-bq$, with $q\in\Z$. There is at least one: choose $q=-|a|$.
If $a\ge0$, then $a+b|a|\ge0$; if $a<0$, then
$a+b|a|=(b-1)|a|\ge0$ because $b\ge1$.
By well ordering, this set has a least element $r=a-bq$.
If $r\ge b$, then $r-b=a-b(q+1)$ is still nonnegative
but is smaller than $r$, a contradiction.
Hence $0\le r<b$, and $a=bq+r$.

\emph{Uniqueness.} Suppose also $a=bq'+r'$ with $0\le r'<b$.
Subtracting the equations gives
\[
b(q-q')=r'-r.
\]
The right side is strictly between $-b$ and $b$.
The only multiple of the positive integer $b$ in that
interval is $0$. Thus $r'=r$, and then $q'=q$.
The sign of $a$ was unrestricted throughout.
\end{proof}

% t1-pass2 why
\paragraph{Why this works.}
The remainder has two jobs, and they are different. Well ordering gives you a smallest nonnegative value among $a-bq$; subtracting another $b$ shows it must be smaller than $b$. That proves existence. For uniqueness, suppose two remainders satisfy the same bounds. Their difference would be a multiple of $b$ with absolute value less than $b$. Only zero can do that. Keep the two jobs separate: a method that finds an answer has not automatically proved it is the only answer.


An \emph{additive subgroup of $\Z$} is a non-empty set $H$ of integers closed under subtraction: $x,y\in H$ implies $x-y\in H$. This automatically gives $0=x-x\in H$, negatives $-x=0-x\in H$, and sums $x+y=x-(-y)\in H$. Conversely a set containing $0$ and closed under addition and negatives is closed under subtraction. Thus the unfamiliar word ``subgroup'' here means precisely a set in which ordinary addition and subtraction never take us outside. We write $d\Z=\{dk:k\in\Z\}$.

\begin{think}
The multiples of $6$ form a subgroup of $(\Z,+)$: sums and differences of multiples of $6$ are multiples of $6$. What about the set $4\Z+6\Z$ of all $4x+6y$? List its elements between $-10$ and $10$. Is it of the form $d\Z$? For which $d$?
\end{think}

\begin{proposition}\label{prop:subgroups}
Every subgroup of $(\Z,+)$ is of the form $d\Z$ for a unique $d\in\N$.
\end{proposition}
\begin{proof}
The subgroup $\{0\}$ is $0\Z$, and no $d\ge1$ generates it. Otherwise let $H\ne\{0\}$ be a subgroup. It contains a non-zero integer $h$, hence also $-h$, so it contains a positive integer. Let $d$ be its least positive element. Then $d\Z\subset H$. For $a\in H$, divide: $a=dq+r$ with $r=a-dq\in H$ and $0\le r<d$, so $r=0$. Uniqueness: $d$ is the least positive element of $d\Z$.
\end{proof}

\begin{definition}
For $a,b\in\Z$ not both zero, $a\Z+b\Z$ is a subgroup of $\Z$, hence equal to $d\Z$ for a unique $d\ge1$: this $d$ is the \emph{greatest common divisor} $\gcd(a,b)$.
\end{definition}

We use the convention $\gcd(0,0)=0$, corresponding to
$0\Z+0\Z=0\Z$; no positive greatest divisor exists in that case.
For $a\ne0$, $\gcd(a,0)=|a|$. The least common multiple is non-negative,
with $\operatorname{lcm}(a,0)=0$. Products of valuation formulas below
are stated only for non-zero inputs, so they never require an undefined
valuation of zero.

This subgroup viewpoint is also called the \emph{ideal} viewpoint: an ideal of a commutative ring is an additive subgroup stable under multiplication by any element of the ring. Every additive subgroup of $\Z$ already has this property, because multiplication by an integer is repeated addition and subtraction. The notation $a\Z+b\Z$ is literal: it is the set of all $ax+by$ with $x,y\in\Z$. It contains $0$ and is closed under differences, so it really is a subgroup.

The name is justified: $d$ divides $a$ and $b$ (they lie in $d\Z$), and any common divisor of $a$ and $b$ divides every element of $a\Z+b\Z$, in particular $d$. Two immediate but fundamental consequences:

\begin{theorem}[B\'ezout]\label{thm:bezout}
There are $u,v\in\Z$ with $au+bv=\gcd(a,b)$. In particular $a$ and $b$ are coprime if and only if $au+bv=1$ for some integers $u,v$.
\end{theorem}
\begin{proof}
$\gcd(a,b)\in d\Z=a\Z+b\Z$. If $au+bv=1$, then $1\in a\Z+b\Z=d\Z$, so $d=1$; the converse is the first statement.
\end{proof}

% t1-pass2 why
\paragraph{Why this works.}
The set $a\Z+b\Z$ is not merely a way to package a theorem. It consists of the integers you can actually make from $a$ and $b$. Its positive generator $d$ is therefore one of those combinations: $d=au+bv$. When $d=1$, you can make $1$, and conversely making $1$ forces the positive generator to be $1$. That is why coprimality and the B\'ezout equation are the same condition seen from two directions.

\begin{theorem}[Gauss' lemma]\label{thm:gauss}
If $a\mid bc$ and $\gcd(a,b)=1$, then $a\mid c$. In particular, if a prime $p$ divides $bc$, it divides $b$ or $c$.
\end{theorem}
\begin{proof}
Write $au+bv=1$. Then $c=acu+bcv$, and $a$ divides both terms. For the second statement, if $p\nmid b$ then $\gcd(p,b)=1$.
\end{proof}

% t1-pass2 why
\paragraph{Why this works.}
Where did coprimality enter the proof? It gave the equation $au+bv=1$. Multiply that identity by $c$ to get $acu+bcv=c$. If $a$ divides $bc$, both terms on the left are multiples of $a$, so $a$ divides $c$. The argument would fail without a B\'ezout identity equal to $1$. For a prime $p$, the alternative is exact: either it divides $b$ or it is coprime to $b$.

\begin{theorem}[fundamental theorem of arithmetic]\label{thm:fta}
Every integer $n\ge2$ is a product of primes, uniquely up to the order of the factors.
\end{theorem}
\begin{proof}
For existence, use strong induction on $n$. If $n$ is prime there is nothing to prove. Otherwise $n=ab$ with $2\le a,b<n$, so the induction hypothesis factors $a$ and $b$ into primes; multiplying the two lists factors $n$. The base case $n=2$ is prime. Uniqueness, by induction on $n$: if $p_1\cdots p_r=q_1\cdots q_s$, then $p_1$ divides the product of the $q_j$, so by Gauss' lemma (iterated) it divides some $q_j$, hence equals it. Cancel and apply the induction hypothesis.
\end{proof}

% t1-pass2 why
\paragraph{Why this works.}
For existence, take a composite $n=ab$ with both factors smaller than $n$ and factor them by strong induction. For uniqueness, match the first prime on one factor list with an equal prime on the other list using Gauss's lemma; cancel and repeat. Ask what would happen if $1$ counted as a prime. You could insert as many copies of $1$ as you wished, so uniqueness would never have a clear meaning. That small convention protects the whole theorem.

For a prime $p$ and $n\ne0$, the \emph{$p$-adic valuation} $v_p(n)$ is the exponent of $p$ in $|n|$. Uniqueness of factorisation gives $v_p(ab)=v_p(a)+v_p(b)$, and
\[
\begin{gathered}\gcd(a,b)=\prod_pp^{\min(v_p(a),v_p(b))},\qq\operatorname{lcm}(a,b)=\prod_pp^{\max(v_p(a),v_p(b))},\\\gcd(a,b)\operatorname{lcm}(a,b)=|ab|\end{gathered}
\]
for $a,b\ne0$. Indeed, for positive $d,n$, uniqueness of factorisation gives $d\mid n$ iff $v_p(d)\le v_p(n)$ for every $p$ (if $n=dk$ then $v_p(n)=v_p(d)+v_p(k)$; conversely $k=\prod p^{v_p(n)-v_p(d)}$ works). Divisibility is unchanged by signs. So the positive common divisors of $a,b$ are the $d$ with $v_p(d)\le\min(v_p(a),v_p(b))$, and the greatest one, which every common divisor divides, is the stated product; likewise for common multiples. The last formula follows from $\min+\max=v_p(a)+v_p(b)$.

The positive divisors of a fixed $n$ are partially ordered by divisibility. For any two of them the greatest lower bound is their gcd and the least upper bound is their lcm, both of which again divide $n$. Figure~\ref{fig:divisor-lattice} draws the case $n=12$. The picture is the reason the formulas above are not two separate inventions: gcd and lcm are the meet and the join of one and the same order.

\begin{figure}[ht]
\centering
\fig{ch01-divisor-lattice}
\caption{The positive divisors of $12$, with an edge when one divides the other and no divisor lies strictly between them. The greatest lower bound of $4$ and $6$ is $\gcd(4,6)=2$; their least upper bound is $\operatorname{lcm}(4,6)=12$.}
\label{fig:divisor-lattice}
\end{figure}

\subsection*{From a definition to a method}
The set $a\Z+b\Z$ contains every integer you can make as $ax+by$. It is worth stopping to see what that means before using the word \emph{ideal}. With $a=4$ and $b=6$, every combination is even. On the other hand, $6-4=2$ is a combination, so every even integer is a combination. Thus $4\Z+6\Z=2\Z$. The smallest positive member of the set is the gcd, and finding it is already a B\'ezout identity.

The word ``greatest'' in greatest common divisor is less useful than its divisibility property:
\[
\begin{gathered}
d\mid a,\quad d\mid b,\\
c\mid a\text{ and }c\mid b\quad\Longrightarrow\quad c\mid d.
\end{gathered}
\]
This property determines the positive $d$ uniquely: two such positive integers divide each other, so they are equal. It is also exactly the property that survives in rings where there is no numerical ordering.

\textbf{A cancellation warning.} From $6\mid2c$ one cannot conclude $6\mid c$: take $c=3$. Gauss' lemma permits cancellation only when the cancelled factor is coprime to the divisor. More generally, if $g=\gcd(a,b)$ with $a\ne0$, then
\[
a\mid bc\quad\Longleftrightarrow\quad \frac ag\mid c.
\]
Indeed, write $a=ga'$, $b=gb'$ with $\gcd(a',b')=1$, divide the divisibility by $g$, and apply Gauss' lemma to $a'\mid b'c$. This is the arithmetic version of reading the small print before cancelling.

% t1-pass2 examples 7efe951b
\begin{examplex}[division with a negative dividend]
Divide $-17$ by $5$: the remainder must lie in $\{0,1,2,3,4\}$, so $-17=5\cdot(-4)+3$, not $5\cdot(-3)-2$. The quotient is $\lfloor -17/5\rfloor=-4$. Uniqueness fails if negative remainders are allowed, which is why the convention is fixed before any algorithm.
\end{examplex}

\begin{examplex}[the ideal generated by $4$ and $6$]
The combinations $4x+6y$ are exactly the multiples of $2$, since $2=6-4$ and every combination is even. Thus $4\Z+6\Z=2\Z$, and $\gcd(4,6)=2$. The same conclusion from valuations: $\min(v_2(4),v_2(6))=1$ and every odd prime has minimum $0$.
\end{examplex}

\section{The Euclidean algorithm}

% t1-pass2 motive a7ecc398
The gcd of $252$ and $198$ was already visible in the staircase above. Now give the procedure its general proof. Every division replaces $(a,b)$ by $(b,r)$ with the same common divisors, and $0\le r<b$. The first fact preserves the answer. The second forces the calculation to stop. Reading the steps backwards produces the coefficients in a B\'ezout identity. We will also count the steps and see why Fibonacci numbers make Euclid work hardest.

First take $a\ge b\ge1$. Negative inputs are replaced by their absolute values; a zero input is immediate, since $\gcd(a,0)=|a|$ for $a\ne0$. Since $\gcd(a,b)=\gcd(b,a-bq)$ (the two pairs generate the same subgroup), we may replace $(a,b)$ by $(b,r)$ where $r$ is the remainder of $a$ divided by $b$. Starting with $r_0=a$, $r_1=b$ and dividing repeatedly,
\[
r_{k-1}=q_kr_k+r_{k+1},\qquad0\le r_{k+1}<r_k,
\]
the remainders decrease strictly, so some $r_{n+1}$ is $0$, and $\gcd(a,b)=r_n$, the last non-zero remainder.

\textbf{Extended version.} Keep track of coefficients with $r_k=au_k+bv_k$: start with $(u_0,v_0)=(1,0)$, $(u_1,v_1)=(0,1)$, and set $u_{k+1}=u_{k-1}-q_ku_k$, $v_{k+1}=v_{k-1}-q_kv_k$. By induction $r_k=au_k+bv_k$ for all $k$ (true for $k=0,1$, and $r_{k+1}=r_{k-1}-q_kr_k$ is matched by the coefficients), so at the end $au_n+bv_n=\gcd(a,b)$.

\begin{example}
For $(a,b)=(240,46)$: $240=5\cdot46+10$, $46=4\cdot10+6$, $10=1\cdot6+4$, $6=1\cdot4+2$, $4=2\cdot2$. So $\gcd=2$, and running the coefficients gives $240\cdot(-9)+46\cdot47=2$.
\end{example}

\subsection*{A complete back-substitution, without a leap}
The remainders above already contain the B\'ezout coefficients; we only have to read them backwards:
\[
\begin{aligned}
2&=6-4=6-(10-6)=2\cdot6-10\\
 &=2(46-4\cdot10)-10=2\cdot46-9\cdot10\\
 &=2\cdot46-9(240-5\cdot46)\\
 &=-9\cdot240+47\cdot46.
\end{aligned}
\]
Each substitution uses an equality from one division, so there is no guesswork. A final substitution into $240u+46v=2$ is a cheap and decisive error check. A sign error should not survive merely because the coefficients look impressively large.

\textbf{What all solutions look like.} Suppose $a,b>0$, $d=\gcd(a,b)$, and $ax_0+by_0=c$. Then every integer solution of $ax+by=c$ is
\[
(x,y)=\left(x_0+\frac bd t,\ y_0-\frac ad t\right),\qquad t\in\Z.
\]
To see completeness, subtract the particular solution and divide by $d$: $(a/d)(x-x_0)=-(b/d)(y-y_0)$. Since $a/d$ and $b/d$ are coprime, Gauss' lemma forces $b/d\mid x-x_0$; substitution then forces the displayed $y$. Conversely each displayed pair works. Existence itself is equivalent to $d\mid c$: necessity follows by divisibility, sufficiency by scaling B\'ezout's identity.

\begin{figure}[ht]
\centering\fig{ch01-bezout-line}
\caption{Integer solutions of $3x+5y=1$. The hollow point $(1,-2/5)$ lies on the real line but is not an integer solution. The filled points $(2,-1)$, $(-3,2)$ and $(7,-4)$ differ by $(5,-3)$. A real line contains infinitely many real points; arithmetic asks which of them land on the integer lattice.}
\label{fig:bezout-line}
\end{figure}

\begin{think}
Find a pair $a>b$ with $b<100$ that makes Euclid's algorithm perform as many divisions as possible. Among such pairs, which has the smallest $a$? Try a few pairs by hand before reading on, and look at the quotients $q_k$ you get: what are they for the slowest pairs?
\end{think}

\begin{figure}[ht]
\centering
\fig{ch01-euclid}
\caption{Euclid's algorithm cuts the largest possible squares from an $a\times b$ rectangle; the number of squares of each size is the quotient $q_k$, and the last square has side $\gcd(a,b)$. On the left, $34\times21$: every quotient except the last is $1$ (the last is $2$), the square sizes are Fibonacci numbers, and the algorithm is as slow as possible (Theorem~\ref{thm:lame}). On the right, $45\times16$: large quotients, few steps.}
\label{fig:euclid}
\end{figure}

\begin{theorem}[Lam\'e]\label{thm:lame}
If the algorithm applied to $a>b\ge1$ performs $n$ divisions, then $b\ge F_{n+1}$ and $a\ge F_{n+2}$, where $(F_k)$ is the Fibonacci sequence.
\end{theorem}
\begin{proof}
The divisions produce $r_0=a>r_1=b>\dots>r_{n-1}>r_n\ge1$ and $r_{n+1}=0$. We show $r_k\ge F_{n+2-k}$ by descending induction on $k$. For $k=n$: $r_n\ge1=F_2$. For $k=n-1$: $r_{n-1}>r_n\ge1$ so $r_{n-1}\ge2=F_3$. Then each $q_k\ge1$, so $r_{k-1}=q_kr_k+r_{k+1}\ge r_k+r_{k+1}\ge F_{n+2-k}+F_{n+1-k}=F_{n+3-k}$.
\end{proof}

So the worst inputs are consecutive Fibonacci numbers, and since $F_{k}\approx\varphi^k/\sqrt5$, the number of steps is $O(\log b)$ (Exercise~\ref{ex:lame5}).

\subsection*{How fast is ``fast''?}
Why does Euclid's algorithm need so few divisions even for large inputs? Each new remainder is smaller than the previous divisor, but that alone gives only a crude estimate. The worst case occurs when the remainders shrink as slowly as possible, with quotients equal to $1$. Work backwards from the final remainders and you meet consecutive Fibonacci numbers. Lam\'e's argument turns that pattern into a bound on the number of division steps. It does not count every operation a modern computer uses to perform each division.

For consecutive Fibonacci inputs the quotients are $1$ until the last quotient $2$; this gives equality in the lower bounds for $a$ and $b$ (with the usual $F_1=F_2=1$). Large quotients make remainders collapse quickly; a long chain of quotients $1$ keeps them alive. The algorithm's least cooperative customer is the golden ratio.

\begin{reminder}[A three-question exit check]
Why does replacing $(a,b)$ by $(b,a-qb)$ preserve \emph{all} integer combinations, rather than merely their gcd? Why does strict decrease of non-negative remainders guarantee termination? And why is the last non-zero remainder a B\'ezout combination of the original inputs? The invariant, termination, and output certificate are three different parts of the proof; together they explain the whole algorithm.
\end{reminder}

% t1-pass2 examples a7ecc398
\begin{examplex}[a complete run for $1071$ and $462$]
One has $1071=2\cdot 462+147$, $462=3\cdot 147+21$ and $147=7\cdot 21+0$. The gcd is $21$. Back-substitution gives $21=462-3\cdot 147=7\cdot 462-3\cdot 1071$. The check $7\cdot 462-3\cdot 1071=3234-3213=21$ catches a sign error.
\end{examplex}

\begin{examplex}[all solutions, once one is known]
From $1071\cdot(-3)+462\cdot 7=21$, the solutions of $1071x+462y=21$ are $x=-3+22t$ and $y=7-51t$, because $462/21=22$ and $1071/21=51$. The particular solution moves by the other coefficient. This is the line drawn for $3x+5y=1$ in Figure~\ref{fig:bezout-line}, with different numbers.
\end{examplex}

\begin{beauty}
Where this leads. Bézout's identity is the private-key step in the toy RSA exercise of Chapter~\ref{ch:congruences}, and it is the reason every Euclidean ring is a principal ideal domain in Volume~3. The same division produces continued fractions in Chapter~\ref{ch:continued} and the Euclidean algorithm in $\Z[i]$ in Chapter~\ref{ch:gaussian}. Volume~13 starts when the ideal $a\Z+b\Z$ is no longer generated by a single element.
\end{beauty}

\subsection*{Proof workshop: preserve a divisor, shrink a remainder}

Take $a=252$ and $b=198$. The large pair is awkward. The remainder $54$
is smaller than either positive input. The key equivalence is
\[
d\mid252\ \text{and}\ d\mid198
\quad\Longleftrightarrow\quad
d\mid198\ \text{and}\ d\mid54.
\]
For the forward direction, $d$ divides $252-198=54$. For the
reverse direction, $d$ divides $198+54=252$. Nothing was lost in
either direction. Apply exactly the same two implications to
$(198,54)$, $(54,36)$ and $(36,18)$.

The pairs end at $(18,0)$. A positive integer divides $18$ and $0$
exactly when it divides $18$. So the common divisors of the original
pair are the positive divisors of $18$, and their greatest is $18$.
The general algorithm is this argument repeated. The termination
proof is separate: each new nonzero remainder is a nonnegative
integer strictly smaller than the previous divisor.

\begin{figure}[ht]
\centering
\fig{ch01-remainder-staircase}
\caption{The Euclidean algorithm for $252$ and $198$. Each arrow
replaces a pair by its divisor and remainder, preserving the
common divisors. The second entry shrinks until it reaches zero.}
\label{fig:euclid-remainder-staircase}
\end{figure}


\section{Exercises}

\begin{exercise}[$\star$]\label{ex:bez}
Find all $(x,y)\in\Z^2$ with $91x+35y=7$.
\end{exercise}

\begin{exercise}[$\star$]\label{ex:gcdn}
For $n\in\N$, compute $\gcd\left(n^2+1,(n+1)^2+1\right)$.
\end{exercise}

\begin{exercise}[$\star$]\label{ex:zeros}
\emph{Legendre's formula.} Show that $v_p(n!)=\sum_{k\ge1}\left\lfloor\frac n{p^k}\right\rfloor$, and find the number of zeros at the end of $1000!$.
\end{exercise}

\begin{exercise}[$\star\star$]\label{ex:sqrt}
Show that $\sqrt n$ is rational only when $n\in\N$ is a perfect square. More generally, show that if $x\in\Q$ is a root of a monic polynomial with integer coefficients, then $x\in\Z$.
\end{exercise}

\begin{exercise}[$\star\star$]\label{ex:mersgcd}
Show that $\gcd\left(2^a-1,2^b-1\right)=2^{\gcd(a,b)}-1$ for $a,b\ge1$.
\end{exercise}

\begin{exercise}[$\star\star$]\label{ex:lame5}
Deduce from Lam\'e's theorem that the number of divisions is at most $5$ times the number of decimal digits of $b$.
\emph{Use $F_{k+1}\ge\varphi^{k-1}$ and $\varphi^5>10$.}
\end{exercise}

\begin{exercise}[$\star\star$]\label{ex:frob}
\emph{The coin problem.} Let $a,b\ge1$ be coprime. Show that the largest integer that cannot be written $ax+by$ with $x,y\in\N$ is $ab-a-b$.
\end{exercise}

\begin{exercise}[$\star\star\star$]\label{ex:supercat}
\begin{enumerate}[label=(\alph*)]
\item Show that $\lfloor2x\rfloor+\lfloor2y\rfloor\ge\lfloor x\rfloor+\lfloor y\rfloor+\lfloor x+y\rfloor$ for all real $x,y$.
\item Deduce that $\dfrac{(2m)!\,(2n)!}{m!\,n!\,(m+n)!}$ is an integer for all $m,n\in\N$.
\end{enumerate}
\end{exercise}

\begin{exercise}[$\star\star\star$]\label{ex:n2n}
Show that for every integer $n\ge2$, $n$ does not divide $2^n-1$. \emph{First show that $2^{p-1}\equiv1\pmod p$ for an odd prime $p$ by expanding $(1+1)^p$. Then use the smallest prime factor of $n$ and Exercise~\ref{ex:mersgcd}.}
\end{exercise}

\begin{exercise}[$\star\star$]\label{ex:t1-divisibility-lattice}
Let $a,b>0$, $d=\gcd(a,b)$ and $\ell=\operatorname{lcm}(a,b)$. Prove directly, using B\'ezout and Gauss' lemma, that
\[
a\Z+b\Z=d\Z,\qquad a\Z\cap b\Z=\ell\Z.
\]
Find the sum and intersection for $a=12$, $b=18$. Why is the word ``sum'' not the same as the union of the two sets?
\end{exercise}

\begin{exercise}[$\star\star$]\label{ex:t1-binomial-carry}
Let $p$ be prime and $m,n\in\N$. Show that
\[
v_p\binom{m+n}m=
\sum_{j\ge1}\left(\left\lfloor\frac{m+n}{p^j}\right\rfloor-
\left\lfloor\frac m{p^j}\right\rfloor-
\left\lfloor\frac n{p^j}\right\rfloor\right).
\]
Explain why each summand is $0$ or $1$, and why it records whether adding $m$ and $n$ in base $p$ produces a carry across the position $p^j$. Check the conclusion on $\binom{10}4=210$ for $p=2$.
\end{exercise}

\begin{exercise}[$\star$]\label{ex:t1-euclid1071}
Compute $\gcd(1071,462)$ by the Euclidean algorithm and find integers $u,v$ with $1071u+462v=\gcd(1071,462)$.
\end{exercise}

\begin{beauty}
Every result of this chapter comes from one fact: the subgroups of $\Z$ are the $d\Z$. In Volume~3 this becomes ``$\Z$ is a principal ideal domain'', and the same argument works verbatim for polynomials over a field, Gaussian integers, and every Euclidean ring. B\'ezout's coefficients are how RSA computes private keys , and the coin problem with three or more coins, the \emph{Frobenius problem}, has no comparably simple two-coin formula; computing it when the number of coin types is part of the input is NP-hard.
\end{beauty}

\begin{exercise}[$\star$]\label{pb:valuation}
\emph{How divisible is a random integer?} Let $p$ be prime and $n$ uniform in $\{1,\dots,N\}$.
\begin{enumerate}[label=(\alph*)]
\item Show that $P(v_p(n)=k)\to(1-\frac1p)p^{-k}$ for every $k\ge0$, and that $\mathbb E[v_p(n)]\to\frac1{p-1}$.
\item Show that for distinct primes $p_1,\dots,p_r$ the valuations $v_{p_1}(n),\dots,v_{p_r}(n)$ are asymptotically independent.
\item Show that $\mathbb E\bigl[2^{v_2(n)}\bigr]=\frac12\log_2N+O(1)$, while $P\bigl(2^{v_2(n)}\le2\bigr)\to\frac34$. Explain the
paradox.
\end{enumerate}
\end{exercise}

\begin{exercise}[$\star$]\label{pb:parity}
\emph{A random fraction.} Choose $a,b$ at random and write $\frac ab=\frac{a'}{b'}$ in lowest terms.
\begin{enumerate}[label=(\alph*)]
\item Show that $p\mid b'$ with probability $\frac1{p+1}$.
\item Deduce that the three parity patterns odd/odd, odd/even, even/odd of $\frac{a'}{b'}$ are equally likely.
\item Show that the events $p\mid b'$, for distinct primes $p$, are asymptotically independent.
\end{enumerate}
\end{exercise}

\begin{exercise}[$\star$]\label{pb:fibcoprime}
\emph{Coprime Fibonacci numbers.} Let $a,b$ be uniform in $\{1,\dots,N\}$ and $F_n$ the Fibonacci numbers.
\begin{enumerate}[label=(\alph*)]
\item Show that $P(\gcd(F_a,F_b)=1)\to\frac{15}{2\pi^2}=0.7599\ldots$
\item Show that $\mathbb E[\log\gcd(F_a,F_b)]=\frac{6\log\varphi}{\pi^2}\log N+O(1)$, where $\varphi$ is the golden ratio.
\end{enumerate}
\end{exercise}

\begin{exercise}[$\star\star$]\label{pb:orchard}
\emph{Pólya's orchard (1918).} Let $R\ge2$ be an integer. A tree stands at every lattice point of the disc of radius $R$ except
the origin, and each trunk is a disc of radius $r$. You stand at the origin.
\begin{enumerate}[label=(\alph*)]
\item (Minkowski.) Show that a closed rectangle centred at $0$, of area at least $4$, contains a lattice point other than $0$.
\item Show that if $r\ge\frac1R$, every line of sight from the origin meets a trunk.
\item Show that if $r<\frac1{\sqrt{R^2+1}}$, you can see out of the orchard.
\end{enumerate}
\end{exercise}

\begin{figure}[ht]
\centering
\fig{ch06-visible}
\caption{The lattice points of the disc of radius $11$: visible from the origin (filled) or hidden behind a nearer point (rings).
A proportion close to $\frac6{\pi^2}\approx61\%$ are visible. The rust sight line, towards $(11,4)$,
passes between the trees: in Pólya's orchard (Problem~\ref{pb:orchard}) it escapes when the trunks are thin enough.}
\label{fig:visible}
\end{figure}

\section{Hints}

\begin{hint}{ex:bez}
Begin with the ordinary gcd of $91$ and $35$. Euclid's algorithm can do more than calculate that number: read the divisions backwards to obtain a B\'ezout identity. Scale it to make the right side $7$, then subtract one solution from another to discover the complete family.
\end{hint}
\begin{hint}{ex:gcdn}
The two entries look quadratic, but a gcd divides their difference too. Use that difference to remove the quadratic term, then combine once more to leave a constant. Once you have a short list of possible gcds, test which possibilities really occur.
\end{hint}
\begin{hint}{ex:zeros}
Count prime factors in $n!$ in layers. Every multiple of $p$ contributes one factor, every multiple of $p^2$ contributes a second, and so on. How many multiples of $p^k$ lie between $1$ and $n$? For decimal zeros, remember that each zero needs both a $2$ and a $5$.
\end{hint}
\begin{hint}{ex:sqrt}
For square roots, write a rational in lowest terms and compare prime exponents after clearing its denominator. For the polynomial question, try $x=a/b$ with $\gcd(a,b)=1$: once you multiply the monic equation by $b^d$, which term is the only one not visibly divisible by $b$?
\end{hint}
\begin{hint}{ex:mersgcd}
Euclid's algorithm reduces the pair $(a,b)$ by replacing the larger exponent with a difference. Check that the associated numbers $2^a-1$ and $2^b-1$ undergo a gcd-preserving reduction at the same time. Then follow the exponent algorithm until it stops.
\end{hint}
\begin{hint}{ex:lame5}
If there are $n$ divisions, $b\ge F_{n+1}\ge\varphi^{n-1}$. Take logarithms in base $10$.
\end{hint}
\begin{hint}{ex:frob}
Every integer $N$ has a unique representation $N=ax+by$ with $0\le x<b$. It is a non-negative representation exactly when $y\ge0$.
\end{hint}
\begin{hint}{ex:supercat}
For (a), reduce to $0\le x,y<1$ by subtracting integer parts. For (b), compare the $p$-adic valuations of numerator and denominator with Legendre's formula.
\end{hint}
\begin{hint}{ex:n2n}
For $0<j<p$, show $p\mid\binom pj$ using $p!=\binom pj j!(p-j)!$ and coprime cancellation. Expand $(1+1)^p$ and cancel $2$ modulo an odd $p$. If $p\mid2^n-1$, then $p\mid\gcd(2^n-1,2^{p-1}-1)$. What is $\gcd(n,p-1)$ when $p$ is the smallest prime factor of $n$?
\end{hint}

\begin{hint}{ex:t1-divisibility-lattice}
The sum was the definition of $d$. For the intersection, put $a=da'$, $b=db'$ with $a',b'$ coprime and write a common multiple as $ak$; Gauss' lemma forces $b'\mid k$. Test whether $12+18$ lies in the union.
\end{hint}
\begin{hint}{ex:t1-binomial-carry}
Subtract three instances of Legendre's formula. For the carry statement divide $m$ and $n$ by $p^j$ and call their remainders $r,s$: a carry crosses that position precisely when $r+s\ge p^j$.
\end{hint}

\begin{hint}{ex:t1-euclid1071}
The remainders are $147$ and then $21$. Substitute backwards from the last nonzero remainder.
\end{hint}

\begin{hint}{pb:valuation}
(a) $\#\{n\le N:v_p(n)\ge k\}=\lfloor N/p^k\rfloor$. (c) The mean is dominated by the rare $n$ divisible by a large power of $2$.
\end{hint}

\begin{hint}{pb:parity}
(a) $p\mid b'$ exactly when $v_p(b)>v_p(a)$.
\end{hint}

\begin{hint}{pb:fibcoprime}
$\gcd(F_a,F_b)=F_{\gcd(a,b)}$ (Volume~1, Exercise~1.7).
\end{hint}

\begin{hint}{pb:orchard}
(a) Halve the rectangle and look at its translates by $\Z^2$. (b) A rectangle of length $2R$ and width $2r$ around the line
of sight; a lattice point $z$ with $|z|^2<R^2+1$ has $|z|\le R$. (c) Look towards $(R,1)$.
\end{hint}

\section{Solutions}

\begin{solution}{ex:bez}
\idea Start by finding a single solution. The extended Euclidean algorithm gives numbers $u,v$ with $au+bv=d=\gcd(a,b)$. If $d$ divides $c$, scaling the identity gives a solution of $ax+by=c$. To find \emph{all} solutions, subtract this one from an arbitrary one. The difference solves $aX+bY=0$, which is easier because coprimality tells us how the two coordinates must vary.
\writeup Run Euclid's algorithm:
\[
91=2\cdot35+21,\quad35=21+14,\quad
21=14+7,\quad14=2\cdot7.
\]
The last nonzero remainder is $7$. Back-substitute to turn it into a combination of the original coefficients:
\[
7=21-14=2\cdot21-35=2\cdot91-5\cdot35.
\]
So $(x_0,y_0)=(2,-5)$ is one solution of $91x+35y=7$.

To find every solution, subtract this particular one from an arbitrary solution. The difference satisfies
$91(x-2)+35(y+5)=0$, or $13(x-2)=-5(y+5)$. Since $5$ and $13$ are coprime, Gauss's lemma gives $5\mid x-2$. Write $x=2+5k$ and solve for $y$ to obtain $y=-5-13k$. Conversely, substitution verifies every pair. Thus
\[
(x,y)=(2+5k,-5-13k),\qquad k\in\Z,
\]
is the complete family, not merely a collection of examples.
\end{solution}

\begin{solution}{ex:gcdn}
\idea A gcd divides every integer combination of the two expressions. So look for combinations that cancel the largest power of $n$, then the remaining power, until only a constant is left. That gives a short list of possible gcd values. Do not stop there: for each divisor of the constant, check when it actually divides both original expressions.
\writeup Let
$d=\gcd(n^2+1,(n+1)^2+1)$. A common divisor divides their difference, so
\[
d\mid (n+1)^2+1-(n^2+1)=2n+1.
\]
Use this linear expression to eliminate $n$ from the other entry:
\[
4(n^2+1)-(2n+1)(2n-1)=5.
\]
Hence $d\mid5$ and, because $d>0$, only $d=1$ or $d=5$ is possible.

For $d=5$, the condition $5\mid2n+1$ forces $n\equiv2\pmod5$. This is also sufficient: if $n=5k+2$, the two starting integers are
$25k^2+20k+5$ and $25k^2+30k+10$, both divisible by $5$. Therefore
\[
d=\begin{cases}
5,&n\equiv2\pmod5,\\
1,&\text{otherwise}.
\end{cases}
\]
We first narrowed the answer to two possibilities, then checked exactly when the exceptional one occurs.
\end{solution}

\begin{solution}{ex:zeros}
\idea Count the factors of $p$ inside $n!$ rather than trying to write out the enormous factorial. Every multiple of $p$ contributes one factor; every multiple of $p^2$ contributes an extra factor; and so on. This turns a valuation into a sum of floors. For trailing decimal zeros, pair factors $2$ and $5$. There are always enough twos, so the limiting supply is the number of fives.
\writeup The $p$-adic valuation of a factorial is the sum of the valuations of all its factors:
\[
v_p(n!)=\sum_{m=1}^nv_p(m).
\]
Each factor divisible by $p$ contributes at least one, each factor divisible by $p^2$ contributes one more, and so on. Counting in those layers gives
\[
\begin{aligned}
v_p(n!)
&=\sum_{m=1}^n\sum_{k\ge1}\mathbf1_{p^k\mid m}
=\sum_{k\ge1}\#\{m\le n:p^k\mid m\}\\
&=\sum_{k\ge1}\left\lfloor\frac n{p^k}\right\rfloor.
\end{aligned}
\]
All sums are finite because $p^k>n$ eventually.

A trailing decimal zero requires one factor $2$ and one factor $5$. Factors of $2$ are more plentiful, so the number of zeros in $1000!$ is
\[
v_5(1000!)=\left\lfloor\frac{1000}5\right\rfloor
+\left\lfloor\frac{1000}{25}\right\rfloor
+\left\lfloor\frac{1000}{125}\right\rfloor
+\left\lfloor\frac{1000}{625}\right\rfloor
=200+40+8+1=249.
\]
\end{solution}

\begin{solution}{ex:sqrt}
\idea Try prime valuations first. A square has an even exponent at every prime, which makes many proposed square-root equalities impossible. For the polynomial part, put a rational root in lowest terms as $p/q$ and clear denominators. Look at the resulting integer equation modulo $q$. Every term but the leading one contains an extra factor of $q$; coprimality then tells you what $q$ must divide.
\writeup For $n=0$, the square root is the integer $0$. Now suppose $n>0$ and $\sqrt n=p/q$ for integers $p,q$ with $q\ne0$. Clearing denominators gives $nq^2=p^2$. For every prime $\ell$, compare valuations:
\[
v_\ell(n)=2v_\ell(p)-2v_\ell(q),
\]
which is even. Thus every prime exponent in $n$ is even, and $n$ is a perfect square. The converse is immediate.

For the general claim, let $x=p/q$ be a rational root in lowest terms, with $q\ge1$, of a monic polynomial
\[
P(X)=X^d+c_{d-1}X^{d-1}+\cdots+c_0,\qquad c_j\in\Z.
\]
Multiply $P(p/q)=0$ by $q^d$:
\[
p^d=-q\bigl(c_{d-1}p^{d-1}
+c_{d-2}p^{d-2}q+\cdots+c_0q^{d-1}\bigr).
\]
Hence $q\mid p^d$. But $\gcd(p,q)=1$, so every prime factor of $q$ would contradict coprimality. Therefore $q=1$ and $x\in\Z$. Taking $P(X)=X^2-n$ gives another proof of the square-root statement.
\end{solution}

\begin{solution}{ex:mersgcd}
\idea Euclid's algorithm is a sequence of remainder replacements. Look for the same replacement in the numbers $2^a-1$ and $2^b-1$. The identity $2^{a+b}-1=2^a(2^b-1)+(2^a-1)$ shows how subtraction of exponents changes the gcd without changing its value. Once that step is proved, follow Euclid's algorithm for the exponents all the way to $\gcd(a,b)$.
\writeup Assume $a\ge b\ge1$. Factor a difference of powers:
\[
2^a-1=2^{a-b}(2^b-1)+(2^{a-b}-1).
\]
Replacing one gcd entry by a remainder therefore gives
\[
\gcd(2^a-1,2^b-1)
=\gcd(2^{a-b}-1,2^b-1).
\]
The exponents have undergone exactly one step of the subtractive Euclidean algorithm. Repeat, interchanging the entries when necessary. Their gcd remains $\gcd(a,b)$, and the process eventually reaches $(d,0)$ with $d=\gcd(a,b)$. At that point
\[
\gcd(2^d-1,2^0-1)
=\gcd(2^d-1,0)=2^d-1.
\]
Thus $\gcd(2^a-1,2^b-1)=2^{\gcd(a,b)}-1$. The algorithm on exponents has solved the gcd problem for exponential numbers.
\end{solution}

\begin{solution}{ex:lame5}
\idea Lam\'e bounds $b$ below by a Fibonacci number; Fibonacci numbers grow like $\varphi^k$; and $d$ digits means $b<10^d$. Chain the inequalities and compare $\varphi^5$ with $10$.
\writeup Write $\varphi=(1+\sqrt5)/2$. We first establish $F_{k+1}\ge\varphi^{k-1}$ for every $k\ge1$. The inequality holds at $k=1,2$. If it holds at the two preceding indices, the Fibonacci recurrence gives
\[
F_{k+2}=F_{k+1}+F_k
\ge\varphi^{k-1}+\varphi^{k-2}
=\varphi^k,
\]
because $\varphi^2=\varphi+1$. This completes the induction.

If Euclid's algorithm on a pair with smaller positive input $b$ takes $n$ divisions, Lam\'e's bound gives
$b\ge F_{n+1}\ge\varphi^{n-1}$. Suppose $b$ has $d$ decimal digits, so $b<10^d$. Since $\varphi^5>10$,
\[
\varphi^{n-1}\le b<10^d<\varphi^{5d}.
\]
Because $\varphi>1$, this implies $n-1<5d$, hence $n\le5d$. The result counts division steps, not the bit-level cost of each division.
\end{solution}

\begin{solution}{ex:frob}
\idea First understand all integer representations of $m$: they form one family $(x-kb,y+ka)$ (Exercise~\ref{ex:bez}'s method). Among them pick the one with $0\le x<b$; the question ``is there one with $x,y\ge0$?'' is then decided by the sign of its $y$. Compute experimentally with $a=3$, $b=5$ to see that $7=3\cdot5-3-5$ is the last failure, then prove it.
\writeup Begin with a B\'ezout representation of an arbitrary integer $m$:
$ax+by=m$, possible because $\gcd(a,b)=1$. If $(x',y')$ is another representation, subtraction gives
$a(x-x')=b(y'-y)$. By Gauss's lemma, $b\mid x-x'$, and all integer representations are of the form
\[
(x-kb,\ y+ka),\qquad k\in\Z.
\]
Exactly one has $0\le x-kb\le b-1$. Call it the canonical representation of $m$.

In that representation, if $y<0$, making it nonnegative would require increasing it by at least $a$, hence replacing $x$ by $x-b<0$. Thus $m$ is a nonnegative combination $ax+by$ precisely when the canonical $y$ is nonnegative.

For $m=ab-a-b$, the canonical choice is
$m=a(b-1)+b(-1)$; its $y=-1$, so $m$ is not representable. If $m>ab-a-b$ and $m=ax+by$ canonically, then
\[
by=m-ax\ge m-a(b-1)>-b.
\]
Therefore $y>-1$, and since $y$ is an integer it is nonnegative. Every larger integer is representable.

When $a=1$ or $b=1$, every nonnegative integer is representable and the formula yields $-1$ for the largest nonrepresentable integer, interpreted in the usual convention. Otherwise the largest nonrepresentable nonnegative integer is exactly $ab-a-b$.
\end{solution}

\begin{solution}{ex:supercat}
\idea (a) Floors of sums depend only on fractional parts, so reduce to $x,y\in[0,1[$, where each floor is $0$ or $1$ and only two cases remain. (b) To prove a fraction is an integer, show every $p$-adic valuation is $\ge0$; Legendre's formula writes it as a sum of exactly the expression of (a).
\writeup (a) Write $x=a+\alpha$ and $y=b+\beta$, where $a,b\in\Z$ and $0\le\alpha,\beta<1$. The integer parts contribute $2a+2b$ to both sides, so the proposed inequality reduces to
\[
\lfloor2\alpha\rfloor+\lfloor2\beta\rfloor
\ge\lfloor\alpha+\beta\rfloor.
\]
If $\alpha+\beta<1$, the right side is $0$. Otherwise at least one of $\alpha,\beta$ is at least $1/2$, so the left side is at least $1$. Because $\alpha+\beta<2$, the right side is then exactly $1$. This proves the inequality in both cases.

(b) Apply Legendre's formula to the factorials of the rational expression in the question. Its $p$-adic valuation is the sum, over $k\ge1$, of
\[
\begin{aligned}
e_k={}&\left\lfloor\frac{2m}{p^k}\right\rfloor
+\left\lfloor\frac{2n}{p^k}\right\rfloor
-\left\lfloor\frac m{p^k}\right\rfloor
-\left\lfloor\frac n{p^k}\right\rfloor\\
&-\left\lfloor\frac{m+n}{p^k}\right\rfloor.
\end{aligned}
\]
Part (a), with $x=m/p^k$ and $y=n/p^k$, says every $e_k\ge0$. Thus every prime occurs with nonnegative valuation. A positive rational number with that property is an integer. This is an example of a floor inequality turning into divisibility.
\end{solution}

\begin{solution}{ex:n2n}
\idea A short binomial argument supplies the prime-exponent congruence without borrowing a later theorem. Then suppose $n\mid2^n-1$ and look at one prime $p\mid n$: then $2^n\equiv1$ and also $2^{p-1}\equiv1\pmod p$ (Fermat), so $p$ divides $2^{\gcd(n,p-1)}-1$ by Exercise~\ref{ex:mersgcd}. To get a contradiction we want $\gcd(n,p-1)=1$, which is guaranteed if $p$ is the \emph{smallest} prime factor of $n$: all prime factors of $n$ then exceed $p-1$.
\writeup Suppose, toward a contradiction, that $n\mid2^n-1$ for some integer $n\ge2$. The right side is odd, so $n$ is odd. Let $p$ be its smallest prime factor; then $p$ is odd and $p\mid2^n-1$.

We can obtain Fermat's little theorem from the binomial expansion. For $1\le j<p$, the prime $p$ divides $\binom pj$: in $p!=\binom pj j!(p-j)!$, neither factorial on the right contains $p$, so Gauss's lemma forces $p\mid\binom pj$. Consequently
\[
2^p=(1+1)^p\equiv2\pmod p.
\]
Since $p$ is odd, $2$ is invertible modulo $p$, giving $2^{p-1}\equiv1\pmod p$.

The prime $p$ now divides both $2^n-1$ and $2^{p-1}-1$. By Exercise~\ref{ex:mersgcd},
\[
p\mid2^{\gcd(n,p-1)}-1.
\]
But every prime divisor of $n$ is at least $p$, whereas every prime divisor of $p-1$ is less than $p$. Hence $\gcd(n,p-1)=1$. The last divisibility becomes $p\mid2-1=1$, impossible. Thus $n\nmid2^n-1$ for all $n\ge2$.
\end{solution}

\begin{solution}{ex:t1-divisibility-lattice}
\idea A sum records what linear combinations can make; an intersection records simultaneous divisibility. Divide out the gcd so that Gauss' lemma can do the cancellation.
\writeup By the definition of the gcd and B\'ezout's identity, all integer combinations of $a,b$ form the subgroup
\[
a\Z+b\Z=d\Z,\qquad d=\gcd(a,b).
\]
For the intersection, write $a=da'$, $b=db'$ with $\gcd(a',b')=1$. If $c\in a\Z\cap b\Z$, then $c=da'k$ for some integer $k$, and $db'\mid da'k$. Hence $b'\mid a'k$. Gauss's lemma gives $b'\mid k$, so $da'b'\mid c$. Conversely $da'b'$ is divisible by both $a$ and $b$. Therefore
\[
a\Z\cap b\Z=da'b'\Z=\operatorname{lcm}(a,b)\Z.
\]
This also proves $\operatorname{lcm}(a,b)=ab/d$ without prime valuations.

For $a=12$, $b=18$, the sum is $6\Z$ and the intersection is $36\Z$. Do not confuse the sum with the union: $30=12+18$ lies in $12\Z+18\Z$, but it belongs to neither $12\Z$ nor $18\Z$, so it is absent from their union.
\end{solution}

\begin{solution}{ex:t1-binomial-carry}
\idea Legendre's floors measure how much divisibility survives after the denominator factorials cancel. A floor increases by one precisely when two remainders overflow their place value.
\writeup The binomial coefficient $\binom{m+n}{m}$ is an integer because it counts choices of $m$ objects from $m+n$. Its valuation, by Legendre's formula, is
\[
v_p\binom{m+n}{m}
=\sum_{j\ge1}\left(
\left\lfloor\frac{m+n}{p^j}\right\rfloor-
\left\lfloor\frac m{p^j}\right\rfloor-
\left\lfloor\frac n{p^j}\right\rfloor\right).
\]
To interpret one summand, write $m=Ap^j+r$, $n=Bp^j+s$ with $0\le r,s<p^j$. All large multiples cancel, leaving $\lfloor(r+s)/p^j\rfloor$. This is $1$ when the sum of the lower $j$ base-$p$ digits crosses the boundary at $p^j$, and $0$ otherwise. It includes any carry already made inside those lower digits. Each summand therefore records one carry, and their sum counts all carries.

For a check in base $2$, $4=100_2$ and $6=110_2$ add to $1010_2$ with one carry into the $2^3$ position. Accordingly $\binom{10}4=210$ has $v_2(210)=1$.
\end{solution}

\begin{solution}{ex:t1-euclid1071}
\idea Each remainder is an explicit combination of the previous two, hence of the original inputs.
\writeup Divide until the remainder vanishes:
\[
1071=2\cdot462+147,\qquad
462=3\cdot147+21,\qquad
147=7\cdot21.
\]
The last nonzero remainder is $21$, so $\gcd(1071,462)=21$. Read the second equation backwards:
\[
21=462-3\cdot147
=462-3(1071-2\cdot462)
=7\cdot462-3\cdot1071.
\]
Thus one B\'ezout pair is $(u,v)=(-3,7)$. Substitution confirms $1071(-3)+462(7)=21$.
\end{solution}

\begin{solution}{pb:valuation}
\idea Divisibility by $p^k$ happens exactly once every $p^k$ integers, so the valuation is geometric in the limit. Distinct
primes do not interact, by the Chinese remainder theorem.

\writeup (a) The number of $n\le N$ with $p^k\mid n$ is $\lfloor N/p^k\rfloor=N/p^k+O(1)$, so
$P(v_p(n)=k)=\frac1N(\lfloor N/p^k\rfloor-\lfloor N/p^{k+1}\rfloor)\to p^{-k}-p^{-k-1}$. For the mean,
$\mathbb E[v_p(n)]=\sum_{k\ge1}P(v_p(n)\ge k)=\frac1N\sum_{k\ge1}\lfloor N/p^k\rfloor$, which lies between
$\sum_{k\le\log_pN}(p^{-k}-\frac1N)$ and $\sum_kp^{-k}=\frac1{p-1}$, and tends to $\frac1{p-1}$.

(b) For integers $k_i\ge0$, the condition $p_i^{k_i}\mid n$ for all $i$ is $m\mid n$ with $m=\prod p_i^{k_i}$, which has
probability $\frac1N\lfloor N/m\rfloor\to\frac1m=\prod_ip_i^{-k_i}$. So the joint tail probabilities $P(v_{p_i}(n)\ge k_i\ \forall i)$
converge to products, and by inclusion--exclusion so do the probabilities $P(v_{p_i}(n)=k_i\ \forall i)$.

(c) $2^{v_2(n)}\le2$ means $4\nmid n$, of probability $\to\frac34$. For the mean, the number of $n\le N$ with $v_2(n)=k$ is
$\frac N{2^{k+1}}+\theta_k$ with $|\theta_k|\le1$, and it is $0$ for $2^k>N$. So
\[
\sum_{n\le N}2^{v_2(n)}=\sum_{k\le\log_2N}2^k\Bigl(\frac N{2^{k+1}}+\theta_k\Bigr)=\frac N2\bigl(\lfloor\log_2N\rfloor+1\bigr)+O(N),
\]
since $\sum_{k\le\log_2N}2^k\le2N$. Divide by $N$. There is no paradox: the mean is carried by rare values. The $n$ with
$v_2(n)=k$ have probability about $2^{-k-1}$ and contribute $2^k\cdot2^{-k-1}=\frac12$ each to the mean, and there are
$\log_2N$ such levels.

\tools Volume~2, Chapter~1 (valuations); Volume~3, Chapter~2 (Chinese remainders).
\leads In the language of Volume~13, Chapter~2, a random integer is a sample of Haar measure on $\Z_p$, simultaneously for all $p$:
that is the meaning of (b). The quantity $2^{v_2(n)}$ is the ``largest power of $2$ dividing $n$'', whose average governs the
cost of the binary gcd algorithm.
\end{solution}

\begin{solution}{pb:parity}
\idea Reducing $\frac ab$ removes the common power of $p$, so $p$ stays in the denominator exactly when $b$ has more factors $p$
than $a$. The valuations of $a$ and $b$ are independent geometric variables.

\writeup (a) $p\mid b'$ if and only if $v_p(b)>v_p(a)$. By Problem~\ref{pb:valuation}(a) applied to $a$ and $b$ independently, for
each $J$ the probability that $v_p(b)>v_p(a)$ and $v_p(a)<J$ tends to
\[
\sum_{j<J}\Bigl(1-\frac1p\Bigr)p^{-j}\cdot p^{-j-1},
\]
and the event $v_p(a)\ge J$ has probability at most $p^{-J}$. Letting $J\to\infty$, the probability tends to
$\frac{p-1}{p^2}\cdot\frac1{1-p^{-2}}=\frac1{p+1}$.

(b) By (a) and symmetry, $2\mid b'$ and $2\mid a'$ each have probability $\frac13$, and they cannot both happen. So each of the
three patterns odd/odd, odd/even, even/odd has probability $\frac13$.

(c) For distinct primes $p_1,\dots,p_r$, the event ``$v_{p_i}(b)>v_{p_i}(a)$ and $v_{p_i}(a)<J$ for all $i$'' is a union of
conditions on $(a,b)$ modulo $\prod p_i^{J+1}$, and by Problem~\ref{pb:valuation}(b) its probability tends to the product of the
individual ones; letting $J\to\infty$ as in (a) gives $\prod_i\frac1{p_i+1}$.

\tools Problem~\ref{pb:valuation}; Volume~2, Chapter~1.
\leads Among random fractions in lowest terms the parity classes are equidistributed, a fact that reappears in the action of the
modular group $SL_2(\Z)$ on $\mathbb P^1(\mathbb F_2)$, a set of $3$ points permuted transitively. By (c), the reduced denominator
of a random fraction is divisible by $m$ with probability $\prod_{p\mid m}\frac1{p+1}$ for squarefree $m$.
\end{solution}

\begin{solution}{pb:fibcoprime}
\idea Fibonacci numbers inherit their gcds from their indices, and $F_g=1$ only for $g=1,2$.

\writeup (a) By Volume~1, Exercise~1.7, $\gcd(F_a,F_b)=F_{\gcd(a,b)}$. Since $F_1=F_2=1$ and $F_g\ge2$ for $g\ge3$, the gcd is $1$ exactly
when $\gcd(a,b)\in\{1,2\}$, which by Problem~\ref{pb:gcdmean}(a) has probability tending to $\frac6{\pi^2}(1+\frac14)=\frac{15}{2\pi^2}$.

(b) Binet's formula gives $F_g=\frac{\varphi^g-(-1/\varphi)^g}{\sqrt5}$, so $|\log F_g-g\log\varphi+\log\sqrt5|\le C$ for all $g\ge1$,
for some constant $C$. Hence $\mathbb E[\log F_{\gcd(a,b)}]=\log\varphi\cdot\mathbb E[\gcd(a,b)]+O(1)$, and Problem~\ref{pb:gcdmean}(b)
finishes.

\tools Volume~1, Chapter~1 (Fibonacci gcd, Binet's formula); Problem~\ref{pb:gcdmean}.
\leads The same holds for every strong divisibility sequence, for instance $2^n-1$ (Volume~2, Chapter~1), where
$\gcd(2^a-1,2^b-1)=1$ with probability $\frac6{\pi^2}$. Two random Fibonacci numbers are coprime more often than two random integers.
\end{solution}

\begin{solution}{pb:orchard}
\idea If the trunks are thick, a thin rectangle around any line of sight has area at least $4$ and so contains a tree. If they are
thin, a line of sight with an ``irrational-looking'' slope $\frac1R$ stays far from all lattice points.

\writeup (a) Let $S$ be the rectangle and $K=\frac12S$, of area at least $1$. Suppose first that the area of $K$ is $>1$. The pieces
$K\cap(z+[0,1[^2)$, $z\in\Z^2$, translated back by $-z$ into $[0,1[^2$, have total area $>1$, so two of them overlap: there are
$x\ne y$ in $K$ with $x-y\in\Z^2\setminus\{0\}$. Then $x-y=\frac12(2x)+\frac12(-2y)$ is the midpoint of two points of $S$ ($S$ is
symmetric), so it lies in $S$ (convex). If the area of $K$ is exactly $1$, apply this to $(1+\frac1m)S$ for each $m$: each contains
a non-zero lattice point, and these lie in the bounded set $2S$, so one of them, $z$, lies in $(1+\frac1m)S$ for infinitely many $m$;
since $S$ is closed, $z\in S$.

(b) If $r\ge1$, every direction makes an angle at most $\frac\pi4$ with one of $(\pm1,0),(0,\pm1)$, and the tree there, at distance
at most $\sin\frac\pi4<1$ from the line, blocks it. Let $\frac1R\le r<1$, and let $D$ be a ray from $0$. Take coordinates $(u,v)$ with
$u$ along $D$, and the rectangle $|u|\le R$, $|v|\le r$, of area $4rR\ge4$. By (a) it contains a lattice point $z\ne0$; replacing $z$ by
$-z$ we may take $u\ge0$, and $u>0$ because $|v|\le r<1\le|z|$. Then $|z|^2=u^2+v^2\le R^2+r^2<R^2+1$, and $|z|^2$ is an integer, so
$|z|\le R$: there is a tree at $z$, and its distance $|v|$ to the ray is at most $r$. The trunk meets $D$.

(c) Let $D$ be the ray towards $(R,1)$. The distance from a lattice point $(a,b)$ to the line through $0$ and $(R,1)$ is
$\frac{|a-bR|}{\sqrt{R^2+1}}$. If $a-bR\ne0$ it is at least $\frac1{\sqrt{R^2+1}}>r$, so that trunk misses the line. If $a=bR$, then
$(a,b)=b(R,1)$ with $b\ne0$, and $|(a,b)|\ge\sqrt{R^2+1}>R$: no tree stands there. So no trunk meets $D$.

\tools Volume~1, Chapters~9 and~7 (convexity; determinants as areas); Volume~2, Chapter~1.
\leads This is Problem~239 of Part~VIII in Pólya and Szegő's \emph{Problems and Theorems in Analysis}, after Pólya's 1918 paper.
Part~(a) is the simplest case of Minkowski's theorem, the first theorem of the geometry of numbers, which Volume~13 develops. The
a direct computation gives the exact critical radius for small $R$: it equals the lower bound $\frac1{\sqrt{R^2+1}}$ every time, so the sight
line towards $(R,1)$ of Figure~\ref{fig:visible} is the best one.
\end{solution}

\section*{Sources and further reading}

\begin{itemize}
\item G.~H.~Hardy and E.~M.~Wright, \emph{An Introduction to the Theory of Numbers}, 6th ed., revised by D.~R.~Heath-Brown and J.~H.~Silverman, Oxford University Press, 2008, Chapters~1--2 and~10. The classic, still a joy.
\item K.~Ireland and M.~Rosen, \emph{A Classical Introduction to Modern Number Theory}, 2nd ed., Graduate Texts in Mathematics~84, Springer, 1990, Chapter~1. Unique factorisation in $\Z$, $k[x]$ and other principal ideal domains.
\item D.~E.~Knuth, \emph{The Art of Computer Programming}, Vol.~2: \emph{Seminumerical Algorithms}, 3rd ed., Addison-Wesley, 1997, Section~4.5.3. The full analysis of Euclid's algorithm, including its average number of steps.
\item V.~Shoup, \emph{A Computational Introduction to Number Theory and Algebra}, 2nd ed., Cambridge University Press, 2009 (free online). Number theory with algorithms and their complexity.
\item J.~L.~Ram\'irez Alfons\'in, \emph{The Diophantine Frobenius Problem}, Oxford University Press, 2005. Everything about the coin problem of Exercise~\ref{ex:frob}, including its NP-hardness.
\item Euclid, \emph{Elements}, Book~VII, Propositions~1--2 (the algorithm), in T.~L.~Heath's translation, Dover, 1956.
\end{itemize}
